% =============================================================================
% Introducción
% =============================================================================
\section{Introducción}

% Contextualización específica del tema
% TODO: Describir el contexto actual de la IA aplicada al diagnóstico médico.
  
La inteligencia artificial (IA), específicamente la rama de aprendizaje profundo (Deep Learning), está revolucionando actualmente la detección de enfermedades por medio de imágenes en varias ramas de la medicina como la radiología, ortopedia, dermatología y oftalmología, entre otras, en las que se utilizan tecnologías avanzadas de imagen para diagnosticar, pronosticar y tratar enfermedades. Años atrás, los médicos analizaban estas imágenes manualmente o utilizando diagnósticos asistidos por computadora (CAD); sin embargo, no ayudaban lo suficiente y frecuentemente resultaba mejor realizar todo el proceso manual debido a limitaciones en su rendimiento. No obstante, la llegada del Deep Learning cambió por completo la situación con algoritmos como las redes neuronales convolucionales (CNN), que han sido la base de muchas arquitecturas utilizadas para identificar enfermedades y que, en diferentes sectores de la medicina, han rendido de manera comparable o incluso superior a especialistas experimentados . Hoy en día, la FDA (Food and Drug Administration), agencia federal principal responsable de regular dispositivos o tecnologías médicas, ha autorizado ya una cantidad considerable de modelos para su utilización en hospitales para que funcionen como copilotos bajo supervisión de especialistas.
% Justificación técnica
% TODO: Explicar por qué es relevante estudiar este tema desde el punto de vista técnico. 

 Las CNN llamaron mucho la atención cuando en 2012 la red convolucional AlexNet superó totalmente a los algoritmos rivales en la competición anual de ImageNet. Estos algoritmos avanzados, capaces de capturar patrones complejos en imágenes para detectar patologías en el cuerpo humano, han sido utilizados para crear arquitecturas con propósitos específicos dentro de la medicina como lo son U-Net, 3D U-Net, V-Net, Attention U-Net y GoogLeNet, entre muchas otras. 
 
Las CNN han ocasionado un gran avance en el ámbito de los CAD, demostrando un excelente rendimiento. El entendimiento de cómo están compuestas estas CNN y su respectivo impacto en la automatización de tareas complejas es de suma importancia para la creación de mejores soluciones que ayuden a evitar, mitigar o facilitar el diagnóstico de enfermedades. Técnicamente, las CNN eliminan la necesidad de extraer manualmente características (handcrafted features), un proceso que no era óptimo y dependía en gran medida de la experiencia humana.
% Relevancia científica o industrial
% TODO: Mencionar la importancia en el ámbito de la salud, investigación clínica, etc.



% Objetivo del documento
% TODO: Indicar qué se busca lograr con este documento de investigación.
En esta investigación se busca transmitir conocimientos acerca del Deep Learning y su utilización en el ámbito médico con evidencia científica probada; también se responderá a preguntas sobre cómo ha impactado el flujo de trabajo de profesionales en el sector de la radiología específicamente, cómo medimos la eficiencia de estos algoritmos avanzados, cuáles arquitecturas han sentado las bases para la mayoría de CNN y, por supuesto, revisaremos los desafíos y aspectos éticos que deben ser considerados para implementar esta tecnología en la práctica clínica.
